\documentclass{article}
\usepackage{amsmath, amsthm}
\newcommand{\R}{\mathbb{R}}
\newcommand{\norm}[1]{\left\lVert #1 \right\rVert}
\newcommand{\inner}[2]{\langle #1, #2 \rangle}
\DeclareMathOperator*{\argmax}{arg\,max}
\DeclareMathOperator{\Tr}{Tr}
\def\eps{\varepsilon}
\newtheorem{theorem}{Theorem}
\title{Notes on Optimization}
\begin{document}
\maketitle
\section{Setup}
Let $f \colon \R^n \to \R$ be differentiable. We write $\norm{x}_2$ for the Euclidean norm and $\inner{x}{y}$ for the inner product.\footnote{See \cite{boyd} for details.}
\begin{theorem}[Descent]
If $\norm{\nabla f(x)} \le \eps$ then $x$ is an $\eps$-stationary point, i.e.\ $\argmax_{y} \inner{\nabla f(x)}{y} \le \eps$.
\end{theorem}
\begin{proof}
By \eqref{eq:main} we have
\begin{align}
  f(x + h) &= f(x) + \inner{\nabla f(x)}{h} + O(\norm{h}^2), \label{eq:main}\\
  \Tr(A^\top A) &= \sum_{i,j} A_{ij}^2 .
\end{align}
\end{proof}
\begin{figure}
\centering
\includegraphics[width=0.5\textwidth]{plot.pdf}
\caption{Loss versus iteration for \emph{three} step sizes.}
\end{figure}
\begin{itemize}
\item First \textbf{point} with $x_i^{(k)}$.
\item Second point, see \url{https://example.com}.
\end{itemize}
\begin{verbatim}
for i in range(10): print(i % 3)
\end{verbatim}
\end{document}
